\providecommand{\rank}{\operatorname{r}}
\providecommand{\RR}{\mathbb R}
\providecommand{\CC}{\mathbb C}
\providecommand{\KK}{\mathbb K}

% ---------- Exercise 65 ----------
\setcounter{enumi}{64}
\item {\bfseries\boldmath
Consider the real vector space $\mathcal F(\RR)$ of all real-valued functions defined on $\RR$. State, with justification, whether each of the following sets is a vector subspace of $\mathcal F(\RR)$.
}

\begin{enumerate}[label=\alph*)]

\item $F_1=\{f\in\mathcal F(\RR):\forall x\in\RR,\ f(-x)=-f(x)\}$ (the set of odd functions).

\begin{answer}
All operations in $\mathcal F(\RR)$ are pointwise.

Yes. The zero function is odd. If $f,g$ are odd, then
\[
 (\alpha f+\beta g)(-x)
 =-\alpha f(x)-\beta g(x)=-(\alpha f+\beta g)(x).
\]
\end{answer}

\item $F_2=\{f\in\mathcal F(\RR):\forall x\in\RR,\ f(x)=f(x+2\pi)\}$ (the set of periodic functions with period $\pi$).

\begin{answer}
Yes. The defining condition is $f(x+2\pi)=f(x)$ for every $x$.
It holds for the zero function and is preserved by linear combinations:
\[
 (\alpha f+\beta g)(x+2\pi)=\alpha f(x)+\beta g(x).
\]
This is $2\pi$-periodicity; the parenthetical reference to period $\pi$ in
the sheet does not agree with its displayed equation.
\end{answer}

\item $F_3=\{f\in\mathcal F(\RR):f\text{ is differentiable}\}$.

\begin{answer}
Yes. The zero function is differentiable, and linear combinations of
differentiable functions are differentiable, with
$(\alpha f+\beta g)'=\alpha f'+\beta g'$.
\end{answer}

\item $F_4=\{f\in\mathcal F(\RR):f\text{ has a second derivative and }\forall x\in\RR,\ f^{\prime\prime}(x)-3f^{\prime}(x)+2f(x)=0\}$.

Note: for any $a,b\in\RR$, the function $f:\RR\to\RR$ given by $f(x)=ae^x+be^{2x}$ belongs to $F_4$ (verify this).

\begin{answer}
Yes. Put $L(f)=f''-3f'+2f$. For twice differentiable $f,g$,
$L(\alpha f+\beta g)=\alpha L(f)+\beta L(g)$, and $L(0)=0$.
Therefore the solutions of $L(f)=0$ form a subspace.
For the additional check, if $f=ae^x+be^{2x}$, then
\[
 f'=ae^x+2be^{2x},\qquad f''=ae^x+4be^{2x},
\]
so $f''-3f'+2f=(1-3+2)ae^x+(4-6+2)be^{2x}=0$.
\end{answer}

\item $F_5=\{f\in\mathcal F(\RR):f\text{ has a second derivative and }\forall x\in\RR,\ f^{\prime\prime}(x)-f^{\prime}(x)+f(x)=\sin x\}$.

Note: the function $f:\RR\to\RR$ given by $f(x)=\cos x$ belongs to $F_5$ (verify this).

\begin{answer}
No. The zero function does not satisfy
$f''-f'+f=\sin x$ for all $x$. The stated example does satisfy it:
for $f=\cos x$, $f'=-\sin x$, $f''=-\cos x$, and
$f''-f'+f=\sin x$.
\end{answer}
\end{enumerate}
